\chapter{تفاعل فيتيغ وتفاعلات أخرى مكوِّنة للرابطة C=C}\label{ch:b2:wittig}

تجذب أنثى الذبابة المنزلية الذكور بهيدروكربون من 23 ذرة كربون فيه رابطة مزدوجة واحدة، بين
الكربونين 9 و10، بهندسة (\textit{Z}). ولصنعه بكميات كبيرة، يجب على الكيميائي أن يضع تلك الرابطة
المزدوجة الوحيدة هناك بالضبط، وبتلك الهندسة: فالحذف كان سيبعثرها على عدة
مواضع ويعطي الهندستين كلتيهما. ويبني هذا الفصل الروابط المزدوجة بالطريقة الأخرى، بوصل
قطعتين عند مجموعة كربونيل، بحيث تظهر الرابطة المزدوجة حيث كانت \ce{C=O}.
% ledger: use:muscalure

\begin{recall}
\Cref{ch:b2:enolates-aldol}: \omterm{def:b2:enolates-aldol:aldol}{التكاثف الألدولي}. و\cref{ch:b2:alkene-redox}: هدرجة
الألكاين إلى ألكين (\textit{Z}) على حفاز مسموم. ومجلد السنة الأولى: الحذوف E1 وE2 وقاعدة زايتسيف،
والاستبدال الثنائي الجزيئية، وكواشف غرينيار وانقلاب القطبية، 
والواصفان (\textit{Z})/(\textit{E}). والمجلد المدرسي: الاقتصاد الذري.
\end{recall}

\section{أملاح الفوسفونيوم والإيليدات}

\begin{definition}[إيليدات الفوسفونيوم]\label{def:b2:wittig:ylide}
\emph{ملح الفوسفونيوم}\index{ملح الفوسفونيوم} $\mathrm{R_3P^+{-}CH_2R'\ X^-}$ يُحضَّر باستبدال ثنائي
الجزيئية لألكان مهلجن بفوسفين، عادة ثلاثي فينيل الفوسفين \ce{Ph3P}. ونزع
هيدروجين من الكربون المرتبط بالفوسفور يعطي \emph{إيليد فوسفونيوم}\index{إيليد فوسفونيوم}،
وهو جزيء متعادل فيه شحنتان متعاكستان متجاورتان، $\mathrm{Ph_3P^+{-}CH^-{-}R'}$، ويُكتب أيضًا $\mathrm{Ph_3P{=}CHR'}$. 
و\emph{الإيليد المثبَّت}\index{إيليد مثبت} يحمل على ذلك الكربون مجموعة ساحبة للإلكترونات
(إستر، أو كيتون، أو \omterm{def:b2:amines:nitrile}{نتريل}) تجعل الشحنة السالبة غير متمركزة.
\end{definition}

\begin{omfigure}
\setchemfig{atom sep=1.9em}
\schemestart
\chemfig{Ph_3P} \+ \chemfig{CH_3-I}
\arrow{->[S$_\mathrm{N}$2]}[,0.9]
\chemfig{Ph_3P^{+}-CH_3}\;\chemfig{I^{-}}
\arrow{->[\ce{BuLi}][$-$ \ce{BuH}, \ce{LiI}]}[,1.3]
\chemfig{Ph_3P^{+}-\chemabove{C}{\scriptstyle -}H_2}
\arrow{<->}[,0.8]
\chemfig{Ph_3P=CH_2}
\schemestop
\par\vspace{6pt}

{\small ميثيليدين ثلاثي فينيل الفوسفوران، أبسط الإيليدات. يزيح ثلاثي فينيل الفوسفين اليوديد
من اليودوميثان؛ ثم ينزع البيوتيل ليثيوم بروتونًا من مجموعة الميثيل. والبنيتان على اليمين
طريقتان لكتابة الجزيء نفسه: والمفصولة الشحنتين هي الأقرب إلى الواقع.}
\end{omfigure}

\begin{proposition}[تحضير الإيليد]\label{prop:b2:wittig:ylide-acidity}
\omterm{def:b2:wittig:ylide}{ملح الفوسفونيوم} أكثر حموضة بكثير عند الكربون المجاور للفوسفور من الألكان. ومع ذلك يحتاج ملح فوسفونيوم
ألكيلي بسيط إلى قاعدة قوية جدًا (البيوتيل ليثيوم، أو هيدريد الصوديوم، أو أميد الصوديوم) في ظروف جافة
خاملة؛ أما الملح الذي تحمل مجموعته \ce{CH2} أيضًا مجموعة كربونيل فيُنزع بروتونه بقاعدة ضعيفة
(كربونات الصوديوم، هيدروكسيد الصوديوم)، ويمكن عزل \omterm{def:b2:wittig:ylide}{الإيليد المثبت} المتكوّن وتخزينه.
\end{proposition}

\begin{proof}[تعليل]
الفوسفور الموجب المجاور للأنيون الكربوني يثبّته بشحنته وبقابلية الاستقطاب لذرة
الفوسفور الكبيرة: فالإيليد البسيط أنيون كربوني مثبت، لكنه أقل تثبيتًا من \omterm{def:b2:enolates-aldol:enolate}{أيون الإينولات}،
ومن هنا الحاجة إلى القاعدة القوية. والكربونيل المجاور ينشر الشحنة السالبة على الأكسجين، كما في \omterm{def:b2:enolates-aldol:enolate}{أيون الإينولات}
(\cref{prop:b2:enolates-aldol:alpha-acidity})، إضافة إلى الفوسفور: فيصير الحمض المرافق
حمضيًا بما يكفي لقاعدة ضعيفة، ولا يعود الإيليد، الأقل قاعدية بكثير، يتفاعل مع الماء أو الهواء.
\end{proof}

\section{تفاعل فيتيغ}

\begin{definition}[تفاعل فيتيغ]\label{def:b2:wittig:wittig}
\emph{تفاعل فيتيغ}\index{تفاعل فيتيغ} \omterm{def:b2:wittig:ylide}{لإيليد فوسفونيوم} مع ألدهيد أو كيتون
يعطي ألكينًا وأكسيد ثلاثي فينيل الفوسفين:
\begin{equation*}
\mathrm{Ph_3P{=}CHR + R'CHO \longrightarrow R'CH{=}CHR + Ph_3P{=}O}.
\end{equation*}
ويمر عبر \emph{أوكسافوسفيتان}\index{أوكسافوسفيتان}، وهو حلقة رباعية من الفوسفور وكربونين
والأكسجين، يتكوّن بارتباط كربون الإيليد بكربون الكربونيل بينما يرتبط أكسجين الكربونيل
بالفوسفور.
\end{definition}

\begin{omfigure}
\setchemfig{atom sep=1.9em}
\schemestart
\chemfig{Ph_3P=CH-R} \+ \chemfig{O=CH-R'}
\arrow{->[إضافة]}[,1.0]
\chemfig{Ph_3P?-[6]CH(-[4]R)-[0]CH(-[0]R')-[2]O?}
\schemestop

\medskip
\schemestart
\arrow{->[حذف]}[,1.1]
\chemfig{R-CH=CH-R'} \+ \chemfig{Ph_3P=O}
\schemestop
\par\vspace{6pt}

{\small \omterm{def:b2:wittig:wittig}{تفاعل فيتيغ}. يرتبط كربون الإيليد وكربون الكربونيل معًا بينما يرتبط الأكسجين
بالفوسفور، في خطوة واحدة: \omterm{def:b2:wittig:wittig}{الأوكسافوسفيتان}. ثم تنفتح حلقته في الاتجاه الآخر، فتنكسر
الرابطتان \ce{P-C} و\ce{C-O}: وتتكوّن الرابطة \ce{C=C} بين الكربونين المتفاعلين السابقين، 
وتتكوّن الرابطة \ce{P=O} لأكسيد ثلاثي فينيل الفوسفين في الوقت نفسه.}
\end{omfigure}

\begin{proposition}[رابطة مزدوجة في مكانها بالضبط]\label{prop:b2:wittig:regiospecific}
في \omterm{def:b2:wittig:wittig}{تفاعل فيتيغ} تصل الرابطة \ce{C=C} الجديدة كربون الكربونيل السابق بكربون الإيليد السابق،
ولا شيء غير ذلك.
\end{proposition}

\begin{proof}[تعليل]
الروابط الوحيدة المتكوّنة والمنكسرة هي روابط الحلقة الرباعية: كربون الإيليد مع كربون
الكربونيل، ثم الرابطتان \ce{C-P} و\ce{C-O}. ولا ينتقل أي هيدروجين ولا يتكوّن كاتيون كربوني ولا أنيون كربوني
على طول السلسلة، فلا يمكن أن تنزاح الرابطة المزدوجة. أما الحذف، فعلى العكس، يختار بين
الهيدروجينات على عدة كربونات مجاورة (قاعدة زايتسيف)، وقد يُعاد ترتيب تفاعل E1.
\end{proof}

\begin{proposition}[القوة الدافعة]\label{prop:b2:wittig:driving}
تكوّن الرابطة \ce{P=O} لأكسيد ثلاثي فينيل الفوسفين يجعل \omterm{def:b2:wittig:wittig}{تفاعل فيتيغ} مواتيًا بقوة
وغير عكوس.
\end{proposition}

\begin{proof}[تعليل]
يستبدل التفاعل برابطة \ce{C=O} ورابطة \ce{P-C} (وفي الإيليد تكون \ce{P=C} في معظمها رابطة أحادية
\ce{P+-C-}) رابطةَ \ce{C=C} ورابطة \ce{P=O}. والرابطة \ce{P=O} في أكسيد الفوسفين
قوية جدًا، أقوى بكثير من الرابطة \ce{P-C} المفقودة، و\ce{C=O} المفقودة تعوّضها جزئيًا
\ce{C=C} المتكوّنة: فالحصيلة نزول واضح. والفوسفور هو مصرف الأكسجين في
التفاعل.
\end{proof}

\section{الانتقائية الفراغية}

\begin{proposition}[هندسة الألكين]\label{prop:b2:wittig:stereo}
مع الألدهيد، يعطي الإيليد غير المثبت (\ce{Ph3P=CHR}، وR مجموعة ألكيل) الألكين (\textit{Z}) أساسًا،
في ظروف خالية من الأملاح. ويعطي \omterm{def:b2:wittig:ylide}{الإيليد المثبت} (\ce{Ph3P=CH-COOR}، \ce{Ph3P=CH-COR}) الألكين
(\textit{E}) أساسًا. أما \omterm{def:b2:wittig:ylide}{الإيليد المثبت} بمجموعة أريل وحدها فيعطي خلائط.
\end{proposition}

\admitted

هذه مشاهدات. وتفسيرها المعتاد يقرأ \omterm{def:b2:wittig:wittig}{تفاعل فيتيغ} على أنه تنافس بين
السرعات والاستقرارات. فمع إيليد تفاعلي غير مثبت، يتكوّن \omterm{def:b2:wittig:wittig}{الأوكسافوسفيتان} بسرعة 
وعلى نحو غير عكوس، عبر الحالة الانتقالية الأقل ازدحامًا، التي ينتهي فيها البديلان في
الجهة نفسها من الحلقة؛ ثم ينهار مع الاحتفاظ بذلك الترتيب: فيعطي التحكم الحركي
الألكين (\textit{Z}). ومع \omterm{def:b2:wittig:ylide}{إيليد مثبت}، تكون الخطوة الأولى أبطأ وعكوسة؛ فتتحول
الأوكسافوسفيتانات بعضها إلى بعض عبر المواد الأولية، ويفوز الذي يقع بديلاه في
جهتين متقابلتين، الأقل توترًا: الألكين (\textit{E})، تحت التحكم الترموديناميكي.

\begin{omfigure}
\setchemfig{atom sep=1.9em}
\schemestart
\chemfig{Ph_3P=CH-CH_2CH_3} \+ \chemfig{O=CH-Ph}
\arrow{->}[,0.8]
\chemfig{Ph-[:-30]=[:30]-[:90]CH_2CH_3}
\schemestop

\medskip
\schemestart
\chemfig{Ph_3P=CH-COOEt} \+ \chemfig{O=CH-Ph}
\arrow{->}[,0.8]
\chemfig{Ph-[:-30]=[:30]-[:-30]COOEt}
\schemestop
\par\vspace{6pt}

{\small إيليدان، وألدهيد واحد. في الأعلى: يعطي إيليد البروبيليدين غير المثبت أساسًا
(\textit{Z})-1-فينيل بوت-1-ين. وفي الأسفل: يعطي \omterm{def:b2:wittig:ylide}{الإيليد المثبت} أساسًا
(\textit{E})-3-فينيل بروبينوات الإيثيل (سينامات الإيثيل). ولا يظهر أكسيد ثلاثي فينيل الفوسفين، المتكوّن في الحالتين،
في الرسم.}
\end{omfigure}

\section{تفاعل هورنر--وادزورث--إيمونز}

\begin{definition}[الفوسفونات وتفاعل HWE]\label{def:b2:wittig:hwe}
\emph{الفوسفونات}\index{فوسفونات} إستر لحمض فوسفوني، $\mathrm{(RO)_2P({=}O){-}R'}$، فيه
رابطة \ce{P-C}. ويُحضَّر بتفاعل ميخائيليس--أربوزوف لفوسفيت ثلاثي الألكيل مع
ألكان مهلجن، $\mathrm{P(OEt)_3 + R'Br \to (EtO)_2P(O)R' + EtBr}$. وفي
\emph{تفاعل هورنر--وادزورث--إيمونز}\index{تفاعل هورنر--وادزورث--إيمونز} (HWE)، يُضاف أنيون
فوسفونات يحمل مجموعة ساحبة للإلكترونات على مجموعته \ce{CH2} إلى ألدهيد فيعطي
ألكينًا وأنيون فوسفات ثنائي الألكيل.
\end{definition}

\begin{omfigure}
\setchemfig{atom sep=1.8em}
\schemestart
\chemfig{EtO-P(=[2]O)(-[6]OEt)-CH_2-COOEt}
\arrow{->[\ce{NaH}][$-$ \ce{H2}]}[,0.9]
\chemfig{EtO-P(=[2]O)(-[6]OEt)-\chemabove{C}{\scriptstyle -}H-COOEt}
\schemestop

\medskip
\schemestart
\arrow{->[\ce{PhCHO}]}[,0.9]
\chemfig{*6(=-=(-[:30]=[:-30]-[:30](=[:90]O)-[:-30]O-[:30]Et)-=-)}
\+
\chemfig{EtO-P(=[2]O)(-[6]O^{-})-OEt}
\schemestop
\par\vspace{6pt}

{\small \omterm{def:b2:wittig:hwe}{تفاعل هورنر--وادزورث--إيمونز} لفوسفونوأسيتات ثلاثي الإيثيل مع البنزالدهيد. ينزع
هيدريد الصوديوم الهيدروجين الواقع بين \omterm{def:b2:wittig:hwe}{الفوسفونات} والإستر؛ ثم يُضاف الأنيون إلى الألدهيد،
وينهار الوسيط الرباعي كما في \omterm{def:b2:wittig:wittig}{تفاعل فيتيغ}. والناتج هو
سينامات الإيثيل (\textit{E})؛ والناتج الثانوي أنيون فوسفات ثنائي الإيثيل، الذواب في الماء.}
\end{omfigure}

\begin{proposition}[لماذا يحب الكيميائيون تفاعل HWE]\label{prop:b2:wittig:hwe}
يعطي تفاعل HWE الألكين (\textit{E}) أساسًا، وناتجه الفوسفوري الثانوي، وهو ملح ذواب في
الماء، يُزال بغسل مائي بسيط؛ وأنيون \omterm{def:b2:wittig:hwe}{الفوسفونات} أيضًا أكثر محبة للنواة من
\omterm{def:b2:wittig:ylide}{الإيليد المثبت} الموافق ويتفاعل مع الكيتونات كذلك.
\end{proposition}

\begin{proof}[تعليل]
أنيون \omterm{def:b2:wittig:hwe}{الفوسفونات} أنيون كربوني مثبت، فتكون إضافته عكوسة ويجري التفاعل
تحت التحكم الترموديناميكي: (\textit{E})، كما مع الإيليدات المثبتة. وخلافًا \omterm{def:b2:wittig:ylide}{لإيليد الفوسفونيوم}، فهو
يحمل شحنة سالبة كاملة، لا يعوّضها كاتيون مجاور: فهو أكثر تفاعلية. 
والناتج الثانوي، \ce{(EtO)2PO2-}، ملح أيوني، بينما أكسيد ثلاثي فينيل الفوسفين جسم صلب عضوي متعادل
يجب فصله عن الناتج بالكروماتوغرافيا أو بالتبلور.
\end{proof}

\section{اختيار طريقة}

\begin{method}[الفك وفق فيتيغ]\label{met:b2:wittig:wittig-disconnection}
\begin{enumerate}
  \item اقطع الرابطة \ce{C=C} في الهدف؛ فيصير أحد الكربونين كربون كربونيل، والآخر كربون الإيليد
        (ملح فوسفونيوم، من هاليد).
  \item من الطريقتين الممكنتين، فضّل التي يكون هاليدها أوليًا (استبدال ثنائي الجزيئية بواسطة
        \ce{Ph3P}) ومركبها الكربونيلي ألدهيدًا.
  \item اختر الكاشف بحسب الهندسة المطلوبة: إيليد غير مثبت من أجل (\textit{Z})، 
        \omterm{def:b2:wittig:ylide}{وإيليد مثبت} أو \omterm{def:b2:wittig:hwe}{فوسفونات} (HWE) من أجل (\textit{E}).
  \item تحقق من أن لا شيء آخر في الجزيء يتفاعل مع القاعدة المستعملة.
\end{enumerate}
\end{method}

\begin{method}[اختيار تفاعل مكوِّن للرابطة C=C]\label{met:b2:wittig:choose-cc}
يلخص الجدول الآتي الطرائق الممكنة.
\begin{center}
\small
\begin{tabular}{@{}lll@{}}
\hline
الطريقة & موضع \ce{C=C} & الهندسة \\
\hline
الحذف E2 أو E1 & قاعدة زايتسيف، خلائط & (\textit{E}) غالبًا، خلائط \\
\omterm{def:b2:enolates-aldol:aldol}{التكاثف الألدولي} & مترافقة مع \ce{C=O} & (\textit{E}) \\
فيتيغ، إيليد غير مثبت & مكان \ce{C=O} السابقة & (\textit{Z}) \\
فيتيغ، \omterm{def:b2:wittig:ylide}{إيليد مثبت}؛ HWE & مكان \ce{C=O} السابقة & (\textit{E}) \\
ألكاين، ثم حفاز مسموم & مكان الرابطة الثلاثية السابقة & (\textit{Z}) \\
\hline
\end{tabular}
\end{center}
\medskip
لوصل قطعتين متقاربتي الحجم عند رابطة \ce{C=C}، يكون \omterm{def:b2:wittig:wittig}{تفاعل فيتيغ} وأقاربه عادة
الخيار الأول؛ أما تبادل الأولفينات، الذي يبادل بين طرفي ألكينين، فيأتي في مجلد السنة
الثالثة.
\end{method}

\begin{history}[الإيليد الذي صار كاشفًا]
وجد غيورغ فيتيغ، وهو يدرس إيليدات الفوسفور في أوائل الخمسينيات، أن ميثيليدين ثلاثي فينيل الفوسفوران
يحوّل البنزوفينون إلى 1،1-ثنائي فينيل إيثين وأكسيد ثلاثي فينيل الفوسفين؛ ومع تلميذه أولريش
شولكوبف حوّل الملاحظة إلى طريقة عامة سنة 1954. وتبنّتها الصناعة بسرعة، ولا سيما من أجل
الفيتامين A، وتقاسم فيتيغ جائزة نوبل في الكيمياء لسنة 1979.
\end{history}

\begin{safety}{\ghs[5mm]{GHS02}\,\ghs[5mm]{GHS05}\,\ghs[5mm]{GHS07}\,\ghs[5mm]{GHS08}}
محاليل البيوتيل ليثيوم قابلة للاشتعال وأكّالة وتتفاعل بعنف مع الماء؛ وهيدريد الصوديوم
يطلق الهيدروجين مع الماء. وثلاثي فينيل الفوسفين ضار وأكّال للعينين، وفوسفيت ثلاثي
الإيثيل قابل للاشتعال وضار؛ وأكسيد ثلاثي فينيل الفوسفين ضار.
% ledger: ghs:BuLi, ghs:PPh3, ghs:triethyl-phosphite, ghs:Ph3PO
\end{safety}

\section{تمارين}

\begin{exercise}[$\star$]\label{exo:b2:wittig:1}
اكتب تحضير الإيليد \ce{Ph3P=CHCH3} من ثلاثي فينيل الفوسفين وألكان مهلجن.
\end{exercise}

\begin{exercise}[$\star$]\label{exo:b2:wittig:2}
أعطِ نواتج \omterm{def:b2:wittig:wittig}{تفاعل فيتيغ} للبنزالدهيد مع الإيليد المحضَّر من اليودوميثان.
\end{exercise}

\begin{exercise}[$\star$]\label{exo:b2:wittig:3}
صنّف إلى مثبت، أو غير مثبت، أو مثبت بأريل: \ce{Ph3P=CHCOOEt}، \ce{Ph3P=CHCH3}،
\ce{Ph3P=CHCOCH3}، \ce{Ph3P=CHC6H5}. أيها يحتاج إلى البيوتيل ليثيوم؟
\end{exercise}

\begin{exercise}[$\star$]\label{exo:b2:wittig:4}
أعطِ ناتج فوسفونوأسيتات ثلاثي الإيثيل وهيدريد الصوديوم والبروبانال، مع هندسته.
\end{exercise}

\begin{exercise}[$\star\star$]\label{exo:b2:wittig:5}
المطلوب ميثيلين هكسان حلقي. قارن طريق فيتيغ انطلاقًا من الهكسانون الحلقي بنزع الماء المحفَّز بالحمض
من 1-ميثيل هكسان حلقي-1-ول.
\end{exercise}

\begin{exercise}[$\star\star$]\label{exo:b2:wittig:6}
فكّك (\textit{E})-1،2-ثنائي فينيل إيثين (الستيلبين) وفق \omterm{def:b2:wittig:wittig}{تفاعل فيتيغ}. لماذا يكون طريق فيتيغ
خيارًا رديئًا للستيلبين (\textit{E}) النقي، وماذا تستعمل بدلًا منه؟
\end{exercise}

\begin{exercise}[$\star\star$]\label{exo:b2:wittig:7}
يتفاعل البروبانال مع \ce{Ph3P=CHCH3} من جهة ومع \ce{Ph3P=CHCOOEt} من جهة أخرى. أعطِ
النواتج وهندستها الغالبة.
\end{exercise}

\begin{exercise}[$\star\star$]\label{exo:b2:wittig:8}
احسب الاقتصاد الذري لميثلنة الهكسانون الحلقي بواسطة \ce{Ph3P=CH2}.
\end{exercise}

\begin{exercise}[$\star\star$]\label{exo:b2:wittig:9}
اكتب تفاعل ميخائيليس--أربوزوف لفوسفيت ثلاثي الإيثيل مع بروموإيثانوات الإيثيل، في خطوتين:
الاستبدال بالفوسفور، ثم نزع الألكيل بالبروميد.
\end{exercise}

\begin{exercise}[$\star\star\star$]\label{exo:b2:wittig:10}
اقترح تخليقًا للمركب 1،6-ثنائي فينيل هكسا-1،3،5-ترايين انطلاقًا من (\textit{E})-بيوت-2-ينديال
\ce{OHC-CH=CH-CHO} بتفاعلي فيتيغ، وقل أي إيليد وبأي كمية.
\end{exercise}

\begin{exercise}[$\star\star\star$]\label{exo:b2:wittig:11}
يمكن تحضير (\textit{E})-4-فينيل بوت-3-ين-2-ون \omterm{def:b2:enolates-aldol:aldol}{بتكاثف ألدولي} أو \omterm{def:b2:wittig:wittig}{بتفاعل فيتيغ}.
اكتب الطريقين وقارن بينهما.
\end{exercise}

\begin{exercise}[$\star\star\star$]\label{exo:b2:wittig:12}
اقترح طريقين إلى (\textit{Z})-أوكت-4-ين، أحدهما \omterm{def:b2:wittig:wittig}{بتفاعل فيتيغ} والآخر عبر ألكاين،
وقل أي النواتج الثانوية يترك كل منهما.
\end{exercise}

\section{مسألة: فيرومون بتفاعل فيتيغ}

\begin{problem}[{مسألة نهاية الأسبوع --- التخليق الراجع لألكين (Z)، والإيليد غير المثبت 
وانتقائيته، والكلفة من أكسيد ثلاثي فينيل الفوسفين، وبديل الألكاين}]
\label{pb:b2:wittig:1}
الهدف: (\textit{Z})-تريكوس-9-ين، \ce{CH3(CH2)7CH=CH(CH2)12CH3}، فيرومون الذبابة المنزلية المذكور في
المقدمة. والكتل المولية (\unit{g/mol}): C 12.011، H 1.008، O 15.999، P 30.974، Br 79.904.
% ledger: use:muscalure, aw:C, aw:H, aw:O, aw:P, aw:Br, ghs:nonanal, ghs:BuLi

\medskip
\noindent\textbf{الجزء الأول --- التخليق الراجع.}
\begin{enumerate}
  \item اكتب الصيغة الجزيئية للهدف وتحقق من أنه ألكين.
  \item أي رابطة يقطعها الفك وفق فيتيغ؟
  \item أعطِ الزوجين الممكنين (مركب كربونيلي، ملح فوسفونيوم).
  \item من أجل الزوج نونانال + إيليد، أي ألكان مهلجن يصنع \omterm{def:b2:wittig:ylide}{ملح الفوسفونيوم}؟
  \item لماذا يناسب ألكان مهلجن أولي هذه الخطوة؟
  \item اكتب تكوّن \omterm{def:b2:wittig:ylide}{ملح الفوسفونيوم}.
\end{enumerate}

\medskip
\noindent\textbf{الجزء الثاني --- الإيليد والهندسة.}
\begin{enumerate}[resume]
  \item أي قاعدة تنزع البروتون من هذا الملح؟ ولماذا لا يصلح هيدروكسيد الصوديوم؟
  \item هل \omterm{def:b2:wittig:ylide}{الإيليد مثبت}؟ وبأي هندسة يتنبأ ذلك؟
  \item اكتب \omterm{def:b2:wittig:wittig}{الأوكسافوسفيتان} المتكوّن مع النونانال.
  \item لماذا يجب إجراء التفاعل جافًا وتحت غاز خامل؟
  \item لماذا تنتهي الرابطة المزدوجة بين الكربونين 9 و10 بالضبط؟
  \item ماذا يعطي نزع الماء المحفَّز بالحمض من تريكوسان-9-ول بدلًا من ذلك؟
\end{enumerate}

\medskip
\noindent\textbf{الجزء الثالث --- أكسيد ثلاثي فينيل الفوسفين.}
\begin{enumerate}[resume]
  \item اكتب المعادلة الإجمالية انطلاقًا من الإيليد والنونانال.
  \item احسب الكتل المولية للهدف، وللنونانال، ولملح الفوسفونيوم ولأكسيد
        ثلاثي فينيل الفوسفين.
  \item احسب الاقتصاد الذري لخطوة فيتيغ (إيليد + ألدهيد).
  \item كيف يُفصل أكسيد ثلاثي فينيل الفوسفين عن هيدروكربون؟
  \item لماذا يكون تكوّن أكسيد ثلاثي فينيل الفوسفين هو القوة الدافعة؟
  \item لماذا لا يناسب تفاعل HWE هذا الهدف؟
\end{enumerate}

\medskip
\noindent\textbf{الجزء الرابع --- طريق الألكاين.}
\begin{enumerate}[resume]
  \item أي ألكاين يعطي الهدف بالهدرجة، وبأي حفاز؟
  \item ابنِ ذلك الألكاين من دك-1-اين وألكان مهلجن: أي هاليد، وأي قاعدة؟
  \item لماذا يجب أن تتوقف الهدرجة عند الألكين، وكيف يضمن الحفاز ذلك؟
  \item قارن الطريقين: عدد الخطوات، والتحكم في الهندسة، والنواتج الثانوية.
  \item ما الاقتصاد الذري لخطوة الهدرجة؟
  \item اذكر كتلة أكسيد ثلاثي فينيل الفوسفين المنتَجة معًا لكل غرام من الفيرومون بطريق
        فيتيغ.
\end{enumerate}
\end{problem}
